\chapter{可压缩流动}

\intro{当流体的速度与声速可比拟时，密度的变化将不可忽略，流动呈现"可压缩"特征。本章从热力学基础出发，系统建立可压缩流动的理论框架，内容涵盖等熵流动、激波理论以及高超音速流动。}

\section{热力学基础}

\subsection{状态方程与基本热力学量}

在可压缩流动中，流体的热力学状态需要用两个独立的状态参数来描述。对于完全气体，压力$p$、密度$\rho$和温度$T$满足状态方程：

\begin{myformula}
p = \rho R T
\end{myformula}

其中$R$为气体常数。对于空气，$R \approx 287\,\mathrm{J/(kg\cdot K)}$。

流体的内能$e$是分子运动能量的宏观体现。对于量热完全气体，内能仅是温度的函数：

\begin{myformula}
e = c_v T
\end{myformula}

其中$c_v$为定容比热容。类似地，焓$h$定义为：

\begin{myformula}
h = e + \frac{p}{\rho} = e + RT = c_p T
\end{myformula}

其中$c_p$为定压比热容。对于完全气体，存在Mayer关系：

\begin{myformula}
c_p - c_v = R
\end{myformula}

\begin{mydefinition}{比热比}
比热比（绝热指数）定义为定压比热容与定容比热容之比：
\begin{myformula}
\gamma = \frac{c_p}{c_v}
\end{myformula}
对于空气（双原子分子），$\gamma = 1.4$。对于单原子气体，$\gamma = \frac{5}{3} \approx 1.667$。
\end{mydefinition}

利用比热比，可将$c_p$和$c_v$用$R$和$\gamma$表示：

\begin{myformula}
c_v = \frac{R}{\gamma-1},\qquad c_p = \frac{\gamma R}{\gamma-1}
\end{myformula}

\bigskip

\noindent\textbf{熵与热力学第二定律。}熵$s$是表征系统无序程度的状态函数。对于可逆过程，熵的变化满足：

\begin{myformula}
T\,\mathrm{d}s = \mathrm{d}e + p\,\mathrm{d}\!\left(\frac{1}{\rho}\right) = \mathrm{d}h - \frac{1}{\rho}\mathrm{d}p
\end{myformula}

对于量热完全气体，从参考状态$(p_0,\rho_0,T_0)$积分可得：

\begin{myformula}
\frac{s-s_0}{c_v} = \ln\left[\frac{T}{T_0}\left(\frac{\rho_0}{\rho}\right)^{\gamma-1}\right] = \ln\left[\frac{p}{p_0}\left(\frac{\rho_0}{\rho}\right)^{\gamma}\right]
\end{myformula}

\subsection{等熵关系}

如果流动过程是可逆且绝热的（即等熵的），则$\mathrm{d}s = 0$，由此导出等熵关系式：

\begin{myformula}
\frac{p}{\rho^{\gamma}} = \text{const},\qquad \frac{T}{\rho^{\gamma-1}} = \text{const},\qquad \frac{p^{\frac{\gamma-1}{\gamma}}}{T} = \text{const}
\end{myformula}

通常以滞止状态（速度为0的状态）为参考，得到滞止参数与当地参数的关系：

\begin{myformula}
\frac{p_0}{p} = \left(1 + \frac{\gamma-1}{2}Ma^{2}\right)^{\frac{\gamma}{\gamma-1}}
\end{myformula}

\begin{myformula}
\frac{\rho_0}{\rho} = \left(1 + \frac{\gamma-1}{2}Ma^{2}\right)^{\frac{1}{\gamma-1}}
\end{myformula}

\begin{myformula}
\frac{T_0}{T} = 1 + \frac{\gamma-1}{2}Ma^{2}
\end{myformula}

其中下标$0$表示滞止状态，$Ma$为当地Mach数。

\subsection{临界状态}

当流动速度恰好等于当地声速时（$Ma=1$），称流动处于临界状态。临界状态下的参数用上标$*$表示。将$Ma=1$代入等熵关系式，得到：

\begin{myformula}
\frac{T^{*}}{T_0} = \frac{2}{\gamma+1},\qquad \frac{p^{*}}{p_0} = \left(\frac{2}{\gamma+1}\right)^{\frac{\gamma}{\gamma-1}},\qquad \frac{\rho^{*}}{\rho_0} = \left(\frac{2}{\gamma+1}\right)^{\frac{1}{\gamma-1}}
\end{myformula}

对于空气（$\gamma=1.4$），$T^{*}/T_0 = 0.8333$，$p^{*}/p_0 = 0.5283$，$\rho^{*}/\rho_0 = 0.6339$。

\bigskip

\noindent\textbf{临界声速$a^{*}$}在可压缩流动理论中具有特殊地位。它定义为：

\begin{myformula}
a^{*} = \sqrt{\gamma R T^{*}} = \sqrt{\frac{2\gamma}{\gamma+1}RT_0}
\end{myformula}

临界声速$Ma^{*}$（有时记作$M^{*}$）定义为速度与临界声速之比：

\begin{myformula}
Ma^{*} = \frac{U}{a^{*}}
\end{myformula}

$Ma^{*}$与通常的Mach数的关系为：

\begin{myformula}
Ma^{*2} = \frac{(\gamma+1)Ma^{2}}{2+(\gamma-1)Ma^{2}}
\end{myformula}

当$Ma\to\infty$时，$Ma^{*}\to\sqrt{(\gamma+1)/(\gamma-1)}$，对于$\gamma=1.4$，极限值为$\sqrt{6}\approx 2.449$。

\section{声速与Mach数}

\subsection{声速的物理本质}

声速是小扰动在可压缩介质中的传播速度。考虑一个微弱的压力脉冲在静止气体中传播，由连续性方程和动量方程在扰动线性化的条件下可得声速的表达式：

\begin{myformula}
a = \sqrt{\left(\frac{\partial p}{\partial\rho}\right)_{s}}
\end{myformula}

下标$s$表示该偏导数在等熵条件下计算。对于完全气体，利用等熵关系求导即得：

\begin{myTheorem}{声速公式}
对于量热完全气体，声速仅取决于当地温度：
\begin{myformula}
a = \sqrt{\gamma R T}
\end{myformula}
\end{myTheorem}

海平面标准条件下（$T=288.15\,\mathrm{K}$），空气的声速约为$340.3\,\mathrm{m/s}$。在平流层（$T\approx 216.7\,\mathrm{K}$），声速降至约$295.1\,\mathrm{m/s}$。

\begin{remark}
声速反映的是介质传播扰动的"刚度"。对于不可压缩流体，$\partial p/\partial\rho\to\infty$，因此声速为无穷大——任何扰动瞬间传播到整个流场。这正是不可压缩假设的数学本质。
\end{remark}

\subsection{Mach数及其分类}

\begin{mydefinition}{Mach数}
Mach数定义为当地流速与当地声速之比：
\begin{myformula}
Ma = \frac{U}{a}
\end{myformula}
它是可压缩流动中最重要的无量纲参数，表征了流动动能与内能的相对重要性。
\end{mydefinition}

根据Mach数的大小，可将流动分类：

\begin{table}[H]
\centering
\caption{按Mach数的流动分类}
\label{tab:mach_class}
\begin{tabular}{lll}
\toprule
\textbf{流动类型} & \textbf{Mach数范围} & \textbf{主要特征} \\
\midrule
不可压缩流动 & $Ma < 0.3$ & 密度变化可忽略($<5\%$) \\
亚音速流动 & $0.3 < Ma < 0.8$ & 无激波，全场信息传播 \\
跨音速流动 & $0.8 < Ma < 1.2$ & 混合亚/超音速区，激波 \\
超音速流动 & $1.2 < Ma < 5$ & 激波，气动加热显著 \\
高超音速流动 & $Ma > 5$ & 薄激波层，真实气体效应 \\
\bottomrule
\end{tabular}
\end{table}

\begin{example}
一架客机以$Ma=0.85$巡航时，属于跨音速飞行。机翼上表面局部气流会加速至超音速，形成激波并引起波阻。而SR-71黑鸟侦察机以$Ma=3.2$巡航，属于超音速飞行，前三章所述不可压缩理论完全失效。
\end{example}

\subsection{Mach角与影响区域}

在超音速流动中，扰动源发出的微弱压力波只能在Mach锥内传播。Mach角$\mu$定义为：

\begin{myformula}
\mu = \arcsin\left(\frac{1}{Ma}\right)
\end{myformula}

\begin{figure}[H]
\centering
\begin{tikzpicture}[scale=1.2]
  \def\xshift{0}
  \fill[ch09supersonic!20] (\xshift,0) -- ++(60:4) -- ++(-60:4) -- cycle;
  \draw[ch09machline,dashed] (\xshift,0) -- ++(60:4);
  \draw[ch09machline,dashed] (\xshift,0) -- ++(-60:4);
  \draw[->,ch09flowline] (\xshift-1,0) -- (\xshift+0.5,0);
  \draw[->,ch09flowline] (\xshift-1,0.2) -- (\xshift+0.5,0.2);
  \filldraw (\xshift-0.3,0) circle (0.05) node[below] {$Ma>1$};
  \draw[ch09arc] (\xshift,0) ++(30:1.5) arc (30:-30:1.5);
  \node[ch09angle] at (\xshift+1.3,0.35) {$\mu$};
  \node[ch09label] at (\xshift+3,2) {Mach锥};
  \node[ch09label] at (\xshift+3,-1) {影响区};
  \node[ch09label] at (\xshift-3,1.5) {寂静区};
\end{tikzpicture}
\caption{Mach锥与影响区域示意图}
\label{fig:mach_cone}
\end{figure}

当$Ma=1$时，$\mu=90^{\circ}$，Mach锥退化为垂直于来流的平面；当$Ma\to\infty$时，$\mu\to 0^{\circ}$。这就是为什么在高超音速飞行时激波紧贴物面。

\section{可压缩流动的基本方程}

\subsection{守恒形式}

可压缩流动受质量、动量和能量守恒定律支配。对于无黏、无热传导的可压缩流动（Euler方程），守恒形式的控制方程为：

\vspace{4pt}
\noindent\textbf{连续性方程：}

\begin{myformula}
\frac{\partial\rho}{\partial t} + \nabla\cdot(\rho\mathbf{u}) = 0
\end{myformula}

\noindent\textbf{动量方程：}

\begin{myformula}
\frac{\partial(\rho\mathbf{u})}{\partial t} + \nabla\cdot(\rho\mathbf{u}\otimes\mathbf{u}) + \nabla p = 0
\end{myformula}

\noindent\textbf{能量方程：}

\begin{myformula}
\frac{\partial(\rho E)}{\partial t} + \nabla\cdot[(\rho E + p)\mathbf{u}] = 0
\end{myformula}

\vspace{4pt}

其中总能量密度$E = e + \frac{1}{2}|\mathbf{u}|^{2}$，$e$为内能。这三个方程构成含有五个未知量（$\rho,u,v,w,p$）的四个方程的系统，需要用状态方程封闭。

对于一维定常流动，方程组大幅简化。连续性方程退化为：

\begin{myformula}
\rho U A = \text{const}
\end{myformula}

动量方程退化为（从Euler方程的定常一维形式积分）：

\begin{myformula}
\rho U\,\mathrm{d}U + \mathrm{d}p = 0
\end{myformula}

能量方程给出：

\begin{myformula}
h + \frac{U^{2}}{2} = h_0 = \text{const}
\end{myformula}

其中$h_0$为总焓（滞止焓），沿流线保持恒定。

\subsection{原始变量形式}

以$(\rho,u,v,w,p,e)$（或$T$）为因变量的原始变量形式有时更便于分析和数值求解。将守恒形式展开即得：

\begin{myformula}
\frac{\mathrm{D}\rho}{\mathrm{D}t} + \rho\nabla\cdot\mathbf{u} = 0
\end{myformula}

\begin{myformula}
\rho\frac{\mathrm{D}\mathbf{u}}{\mathrm{D}t} = -\nabla p
\end{myformula}

\begin{myformula}
\rho\frac{\mathrm{D}h}{\mathrm{D}t} = \frac{\mathrm{D}p}{\mathrm{D}t}
\end{myformula}

最后一个方程（能量方程）表明，在等熵流动中，滞止焓守恒等价于$h + U^{2}/2 = \text{const}$。

\subsection{速度势方程与全速势方程}

对于无旋的等熵流动，可引入速度势$\phi$使得$\mathbf{u} = \nabla\phi$。从连续性和动量方程消去密度，得到全速势方程：

\begin{myformula}
\left(a^{2} - \phi_x^{2}\right)\phi_{xx} + \left(a^{2} - \phi_y^{2}\right)\phi_{yy} + \left(a^{2} - \phi_z^{2}\right)\phi_{zz} - 2\phi_x\phi_y\phi_{xy} - 2\phi_y\phi_z\phi_{yz} - 2\phi_z\phi_x\phi_{zx} = 0
\end{myformula}

这是一个拟线性偏微分方程。当$Ma<1$时为椭圆型，$Ma>1$时为双曲型——这解释了亚音速和超音速流动在数学性质上的根本差异。

\section{一维等熵流动}

\subsection{变截面管流}

考虑一维定常等熵管流。质量流量$\dot{m} = \rho U A$为常数。对连续性方程取对数微分，有：

\begin{myformula}
\frac{\mathrm{d}\rho}{\rho} + \frac{\mathrm{d}U}{U} + \frac{\mathrm{d}A}{A} = 0
\end{myformula}

联立动量方程$\rho U\mathrm{d}U + \mathrm{d}p = 0$和声速关系$\mathrm{d}p = a^{2}\mathrm{d}\rho$，消去$\mathrm{d}p$和$\mathrm{d}\rho$，得到著名的面积-速度关系：

\begin{myTheorem}{面积-速度关系}
对于一维定常等熵管流，速度的变化与截面积的变化满足：
\begin{myformula}
\frac{\mathrm{d}U}{U} = -\frac{1}{1-Ma^{2}}\frac{\mathrm{d}A}{A}
\end{myformula}
\end{myTheorem}

这个简洁的关系揭示了亚音速和超音速管流行为的本质差异：

\begin{itemize}
  \item 当$Ma<1$时，$1-Ma^{2}>0$，$\mathrm{d}A$与$\mathrm{d}U$异号——收敛通道加速（$\mathrm{d}A<0\to\mathrm{d}U>0$），扩张通道减速。
  \item 当$Ma>1$时，$1-Ma^{2}<0$，$\mathrm{d}A$与$\mathrm{d}U$同号——收敛通道减速，扩张通道加速。
\end{itemize}

\begin{remark}
这一反直觉的行为是超音速流动的特征：气体在扩张通道中反而加速。物理上，这是因为超音速时密度下降的速度超过了面积增加的速度，迫使流速增加以维持质量流量守恒。
\end{remark}

\subsection{Laval喷管}

Laval喷管（又称收缩-扩张喷管）利用上小节所述的截面变化规律，首先在收缩段将气流加速至音速（喉部），然后在扩张段继续加速至超音速。

\begin{figure}[H]
\centering
\begin{tikzpicture}[scale=1.1]
  % Nozzle contour
  \draw[ch09nozzle,fill=ch09nozzlefill!30] 
    (0,2.5) -- (2,2) .. controls (4,1) and (4,1) .. (6,0.8) -- 
    (6,-0.8) .. controls (4,-1) and (4,-1) .. (2,-2) -- (0,-2.5);
  
  % Centerline
  \draw[dashed,gray] (0,0) -- (8.5,0);
  
  % Flow arrows
  \foreach \x/\y in {1.5/1.6,3.5/1.1,5.5/0.9,7.5/1.3} {
    \draw[->,ch09flowline] (\x-\y*0.15,\y) -- (\x+0.5-\y*0.15,\y);
    \draw[->,ch09flowline] (\x-\y*0.15,-\y) -- (\x+0.5-\y*0.15,-\y);
  }
  
  % Regions
  \node[ch09label] at (2,-1.5) {亚音速};
  \node[ch09label] at (4.5,-1.5) {喉部\\$Ma=1$};
  \node[ch09label] at (7.5,-1.5) {超音速};
  
  % Throat marker
  \draw[ch09throat] (4,1) -- (4,-1);
  \node[ch09label] at (4,1.3) {喉部};
  
  \draw[<->,ch09dim] (3,0.5) -- (5,0.5);
  \node[ch09label] at (4,0.75) {$A^{*}$};
\end{tikzpicture}
\caption{Laval喷管：收缩段中气流为亚音速（$Ma<1$），在喉部达到音速（$Ma=1$），扩张段中继续加速至超音速（$Ma>1$）。}
\label{fig:laval_nozzle}
\end{figure}

面积比与Mach数的关系由下式给出：

\begin{myformula}
\left(\frac{A}{A^{*}}\right)^{2} = \frac{1}{Ma^{2}}\left[\frac{2}{\gamma+1}\left(1+\frac{\gamma-1}{2}Ma^{2}\right)\right]^{\frac{\gamma+1}{\gamma-1}}
\end{myformula}

其中$A^{*}$为喉部面积（临界面积）。对于给定的面积比$A/A^{*}$，该方程有两个解：一个亚音速解（$Ma<1$）和一个超音速解（$Ma>1$），对应Laval喷管的两种工作模态。

\subsection{壅塞现象}

\begin{mydefinition}{壅塞}
当喷管喉部达到音速（$Ma=1$）时，质量流量达到最大值，进一步降低背压不会增加流量——这一现象称为壅塞（choking）。
\end{mydefinition}

壅塞质量流量（以喉部参数表示）为：

\begin{myformula}
\dot{m}_{\max} = \rho^{*} a^{*} A^{*} = p_0 A^{*} \sqrt{\frac{\gamma}{RT_0}}\left(\frac{2}{\gamma+1}\right)^{\frac{\gamma+1}{2(\gamma-1)}}
\end{myformula}

壅塞是喷管设计和安全分析中的关键概念。燃气轮机、火箭喷管以及事故情况下的管道破裂泄放计算，均涉及壅塞条件的判断。

\begin{example}
某高压储罐（$p_0=10\,\text{bar}$，$T_0=300\,\text{K}$）通过一收敛-扩张喷管向大气排放。喉部直径为$5\,\text{mm}$，出口面积为喉部的$2.4$倍。计算壅塞条件下的质量流量。

查表：对于$\gamma=1.4$，$A/A^{*}=2.4$对应设计Mach数$Ma=2.4$（超音速）。壅塞流量：
\[
\dot{m} = \frac{10^{6}\times \pi(0.0025)^{2}}{\sqrt{287\times 300}}\times\sqrt{1.4}\times\left(\frac{2}{2.4}\right)^{3} \approx 0.078\,\text{kg/s}
\]
\end{example}

\section{正激波}

\subsection{激波的形成}

在超音速流动中，当气流遇到障碍物或背压升高时，不可能通过连续的压缩波来减速——这些压缩波会汇聚成一个有限强度的间断面，即激波。激波是流动参数发生跃变的极薄层（厚度仅为几个分子平均自由程量级）。

取相对激波静止的控制体，上游（来流）的量用下标$1$标记，下游的量用下标$2$标记。激波前后的流动参数通过守恒定律联系。

\subsection{Rankine-Hugoniot关系}

考虑垂直于激波的一维定常流动。跨激波的守恒方程为：

\vspace{4pt}
\noindent\textbf{质量守恒：}

\begin{myformula}
\rho_1 u_1 = \rho_2 u_2
\end{myformula}

\noindent\textbf{动量守恒：}

\begin{myformula}
p_1 + \rho_1 u_1^{2} = p_2 + \rho_2 u_2^{2}
\end{myformula}

\noindent\textbf{能量守恒：}

\begin{myformula}
h_1 + \frac{u_1^{2}}{2} = h_2 + \frac{u_2^{2}}{2} = h_0
\end{myformula}

\vspace{4pt}

对于量热完全气体（$h=c_pT$），从这三个方程可导出激波前后的参数比值。定义激波强度为来流Mach数$Ma_1$，经代数推导得到Rankine-Hugoniot关系式：

\begin{myformula}
\frac{p_2}{p_1} = 1 + \frac{2\gamma}{\gamma+1}(Ma_1^{2} - 1)
\end{myformula}

\begin{myformula}
\frac{\rho_2}{\rho_1} = \frac{u_1}{u_2} = \frac{(\gamma+1)Ma_1^{2}}{(\gamma-1)Ma_1^{2} + 2}
\end{myformula}

\begin{myformula}
\frac{T_2}{T_1} = \left[1 + \frac{2\gamma}{\gamma+1}(Ma_1^{2}-1)\right]\left[\frac{2+(\gamma-1)Ma_1^{2}}{(\gamma+1)Ma_1^{2}}\right]
\end{myformula}

\begin{myformula}
Ma_2^{2} = \frac{1 + \frac{\gamma-1}{2}Ma_1^{2}}{\gamma Ma_1^{2} - \frac{\gamma-1}{2}}
\end{myformula}

\vspace{4pt}

\begin{figure}[H]
\centering
\begin{tikzpicture}[scale=1.1]
  % Shock plane
  \shade[left color=ch09shock!40, right color=ch09shock!10] (2,-2.5) rectangle (2.1,2.5);
  \draw[ch09shock,very thick] (2,-2.5) -- (2,2.5);
  
  % Upstream flow
  \draw[->,ch09flowline,ultra thick] (0,0) -- (1.8,0) node[midway,above] {$u_1 > a_1$};
  \node[ch09label] at (0.8,1.2) {上游(1)};
  \node[ch09label] at (0.8,-2) {$p_1,\rho_1,T_1$};
  \node[ch09label] at (0.8,-2.5) {$Ma_1 > 1$};
  
  % Downstream flow
  \draw[->,ch09flowline,ultra thick] (2.3,0) -- (4.0,0) node[midway,above] {$u_2 < a_2$};
  \node[ch09label] at (3.2,1.2) {下游(2)};
  \node[ch09label] at (3.2,-2) {$p_2,\rho_2,T_2$};
  \node[ch09label] at (3.2,-2.5) {$Ma_2 < 1$};
  
  % Parameters
  \node[ch09label] at (1,0.5) {$p_2 > p_1$};
  \node[ch09label] at (2.8,-1.5) {$\rho_2 > \rho_1$};
  \node[ch09label] at (2.8,-1.9) {$T_2 > T_1$};
  \node[ch09label] at (1,-0.4) {$s_2 > s_1$};
  
  \node at (2,2.8) {\textbf{激波面}};
\end{tikzpicture}
\caption{正激波前后的参数变化}
\label{fig:normal_shock}
\end{figure}

\begin{remark}
从Rankine-Hugoniot关系可知：(1)激波前的流动必为超音速（$Ma_1>1$），激波后的流动必为亚音速（$Ma_2<1$）；(2)激波是压缩的——压力、密度和温度均升高；(3)熵增加——激波是不可逆过程。
\end{remark}

\subsection{Prandtl关系}

\begin{myTheorem}{Prandtl关系}
对于正激波，激波前后的速度与临界声速之间满足：
\begin{myformula}
u_1 u_2 = a^{*2}
\end{myformula}
或以无量纲形式表示：
\begin{myformula}
Ma_1^{*} \cdot Ma_2^{*} = 1
\end{myformula}
\end{myTheorem}

Prandtl关系揭示了激波的物理本质：如果来流为超音速（$Ma_1^{*}>1$），则激波后必为亚音速（$Ma_2^{*}<1$），且二者互为倒数。这一简洁关系在激波管和喷管流动分析中极为有用。

\begin{proof}
从动量方程和连续性方程出发：
\[
p_2 - p_1 = \rho_1 u_1^{2} - \rho_2 u_2^{2} = \rho_1 u_1(u_1 - u_2)
\]
利用声速表达式和能量方程，代入后整理即得$u_1 u_2 = a^{*2}$。
\end{proof}

\subsection{激波管问题}

激波管是产生和测量激波的基本实验装置。初始时刻隔膜两侧为高压和低压气体，隔膜破裂后形成向右传播的激波和向左传播的稀疏波。

利用Rankine-Hugoniot关系、等熵关系以及接触面两侧速度/压力匹配条件，可求解激波管流动的完整参数分布——这是经典的Riemann问题，其精确解为现代可压缩流动数值格式（Godunov格式等）提供了理论基础。

\section{斜激波与膨胀波}

\subsection{斜激波}

当超音速来流遇到楔形体或压缩拐角时，流动参数并非在法向平面内发生跃变，而是沿倾斜的激波面。斜激波的强度取决于来流Mach数$Ma_1$和激波角$\beta$。

将速度矢量沿激波面分解为切向分量和法向分量。切向速度跨越激波连续，只有法向分量经受与正激波相同的跃变。引入等效正激波Mach数$Ma_{1n} = Ma_1\sin\beta$，则斜激波前后的参数关系可通过将$Ma_{1n}$代入正激波的Rankine-Hugoniot关系得到。

\begin{myformula}
\frac{p_2}{p_1} = 1 + \frac{2\gamma}{\gamma+1}(Ma_1^{2}\sin^{2}\beta - 1)
\end{myformula}

\begin{myformula}
\tan\theta = 2\cot\beta\,\frac{Ma_1^{2}\sin^{2}\beta - 1}{Ma_1^{2}(\gamma + \cos 2\beta) + 2}
\end{myformula}

其中$\theta$为气流偏转角（楔形体半角）。

对于给定的$Ma_1$和$\theta$，上式给出$\beta$的两个解：强激波解（$\beta$较大，下游$Ma_2<1$）和弱激波解（$\beta$较小，下游通常$Ma_2>1$）。实际观察到的通常是弱激波解。

\begin{mydefinition}{激波-膨胀波理论}
对于小偏转角，可使用线性化理论（Ackeret公式）近似计算压力分布：
\begin{myformula}
C_p = \frac{p-p_{\infty}}{\frac{1}{2}\rho_{\infty}U_{\infty}^{2}} \approx \frac{2\theta}{\sqrt{Ma_{\infty}^{2}-1}}
\end{myformula}
压缩面（$\theta>0$）：$C_p>0$，压力升高；膨胀面（$\theta<0$）：$C_p<0$，压力降低。
\end{mydefinition}

\subsection{Prandtl-Meyer膨胀波}

超音速气流绕过凸角时，通过一系列Mach波（膨胀波）完成加速，这一过程是等熵的和连续的（与激波的非等熵、不连续形成对比）。由Prandtl和Meyer首先定量描述。

定义Prandtl-Meyer函数$\nu(Ma)$：

\begin{myformula}
\nu(Ma) = \sqrt{\frac{\gamma+1}{\gamma-1}}\arctan\sqrt{\frac{\gamma-1}{\gamma+1}(Ma^{2}-1)} - \arctan\sqrt{Ma^{2}-1}
\end{myformula}

气流绕过凸角$\theta$时，下游Mach数由下式确定：

\begin{myformula}
\nu(Ma_2) = \nu(Ma_1) + \theta
\end{myformula}

\begin{figure}[H]
\centering
\begin{tikzpicture}[scale=1.0]
  % Wall
  \draw[ch09wall,ultra thick] (0,0) -- (4,0);
  \draw[ch09wall,ultra thick] (4,0) -- (6,-1.5);
  
  % Flow
  \draw[->,ch09flowline,ultra thick] (-0.5,0.5) -- (1.5,0.5);
  \draw[->,ch09flowline,ultra thick] (-0.5,0.2) -- (1.5,0.2);
  
  % Expansion fan
  \foreach \i in {1,2,...,6} {
    \draw[ch09fan] (4,0) -- ({4+4*cos(10-40/5*\i)},{0.5-4*sin(10-40/5*\i)});
  }
  
  % Downstream
  \draw[->,ch09flowline,ultra thick] (5,0.1) -- (7.5,-1.8);
  \draw[->,ch09flowline,ultra thick] (5.2,-0.05) -- (7.7,-1.95);
  
  \node[ch09label] at (2,0.8) {$Ma_1 > 1$};
  \node[ch09label] at (7.3,-1.2) {$Ma_2 > Ma_1$};
  \node[ch09label] at (4.5,0.7) {膨胀波扇};
  \draw[ch09arc] (4,0) +(20:0.5) arc (20:-20:0.5);
  \node[ch09angle] at (4.8,-0.15) {$\theta$};
\end{tikzpicture}
\caption{Prandtl-Meyer膨胀波：超音速气流绕过凸角，通过膨胀波扇加速。流动为等熵过程。}
\label{fig:pm_expansion}
\end{figure}

\begin{example}
$Ma_1=2.0$的气流绕过$\theta=10^{\circ}$的凸角。查Prandtl-Meyer表：$\nu(Ma_1=2.0)=26.38^{\circ}$，故$\nu(Ma_2)=26.38^{\circ}+10^{\circ}=36.38^{\circ}$，反查得$Ma_2\approx 2.38$。
\end{example}

\subsection{激波极线图}

激波极线图是分析斜激波的有力工具。对于给定的$Ma_1$，将所有可能的$(\theta,\beta)$对以及对应的$p_2/p_1$绘制成图即得激波极线。其主要特征为：

\begin{itemize}
  \item 存在最大偏转角$\theta_{\max}$，超过此角度的楔形体将产生脱体激波。
  \item $\theta=0$对应Mach波（$\mu$）或正激波（$90^{\circ}$）。
  \item 弱激波支对应滑动接触的上游，强激波支对应下游。
\end{itemize}

对于膨胀波，Prandtl-Meyer理论给出了等熵加速的全部可能解，没有类似"最大偏转角"的限制——超音速气流可通过多次折射实现几乎任意角度的转向。

\section{高超音速流动}

\subsection{高超音速流动的特征}

当Mach数超过约5时，流动进入高超音速范畴。与普通的超音速流动相比，高超音速具有以下鲜明特征：

\begin{itemize}
  \item \textbf{薄激波层。}激波角$\beta$接近Mach角$\mu=\arcsin(1/Ma)$，激波紧贴物体表面，激波层厚度很小。
  \item \textbf{熵层。}激波层内存在强烈的熵梯度——强烈弯曲的弓形激波后，沿流线的熵各不相等。
  \item \textbf{黏性干扰。}激波层与边界层高度重叠甚至合并，无黏假设不再成立。
  \item \textbf{气动加热。}动能转化为内能，边界层内的温度极高。壁面热流$q_w \propto \rho_{\infty} U_{\infty}^{3}$，对速度极其敏感。
  \item \textbf{真实气体效应。}高温下气体分子发生振动激发、离解、电离乃至辐射，比热容不再为常数。量热完全气体假设失效。
\end{itemize}

\subsection{气动加热与热防护}

气动加热是高超音速飞行器设计中面临的首要挑战。壁面的热流率可通过参考焓方法或工程关联式估算。对于钝头体前缘，Fay-Riddell关联式给出驻点热流：

\begin{myformula}
q_w = 0.94(\rho_e\mu_e)^{0.1}(\rho_w\mu_w)^{0.4}\left[1+(Le^{0.52}-1)\frac{h_D}{h_e}\right](h_e - h_w)\sqrt{\left(\frac{\mathrm{d}u_e}{\mathrm{d}x}\right)_s}
\end{myformula}

其中Le为Lewis数，下标$e$表示边界层外缘，$w$表示壁面，$h_D$为离解焓。

热防护系统（TPS）的设计原则包括：（1）烧蚀式热防护——表层材料蒸发/升华带走热量（阿波罗返回舱）；（2）辐射式热防护——高温表面向外辐射散热（航天飞机）；（3）被动隔热——低热导率隔热瓦。

\subsection{真实气体效应}

\noindent\textbf{热力学非平衡。}在高温激波层中，分子平动温度首先升高，而振动温度和化学温度以有限速率弛豫（Landau-Teller模型）。温度的"多温模型"——平动温度$T_{\text{tr}}$、振动温度$T_{\text{vib}}$和电子温度$T_e$——取代单一温度。

\noindent\textbf{化学非平衡。}温度超过约2500 K时（对空气），氧气开始离解，约4000 K时氮气离解，约9000 K时电离开始。化学反应速率（Arrhenius定律）有限，使得化学非平衡效应显著。Damkohler数$Da = \tau_f/\tau_c$（流动时间与化学反应时间之比）是判断非平衡程度的关键参数。

\begin{myformula}
\frac{\mathrm{D}c_s}{\mathrm{D}t} = \dot{\omega}_s(T,\{c_i\})
\end{myformula}

其中$c_s$为组分$s$的质量分数，$\dot{\omega}_s$为由化学反应动力学模型给出的净生成率。

\noindent\textbf{辐射耦合。}在极高空速（$Ma>10$）再入时，激波层气体的辐射不再可忽略。辐射-对流耦合使得能量方程中增加辐射源项，问题的数学复杂性显著提升——能量方程从微分方程变为积分-微分方程。

\begin{remark}
高超音速的真实气体效应使得"Mach数相似律"失效。在常比热假设下，流动无量纲化后的解仅取决于$Ma_{\infty}$和$\gamma$；但在真实气体条件下，来流密度、速度和特征尺度也独立出现——流动不再是相似可缩放的。
\end{remark}

\subsection{航天器再入气动热力学}

再入大气层是所有高超音速理论集大成的工程问题。航天器以约7.8 km/s（第一宇宙速度）再入时，动能转换为热能的速率高达数十至数百MW/m$^{2}$。

\begin{table}[H]
\centering
\caption{典型再入条件对比}
\label{tab:reentry}
\begin{tabular}{lccc}
\toprule
\textbf{任务} & \textbf{Ma} & \textbf{驻点热流} & \textbf{TPS类型} \\
\midrule
阿波罗返回 & 32 & $\sim$40 MW/m$^{2}$ & AVCOAT烧蚀 \\
航天飞机 & 25 & $\sim$15 MW/m$^{2}$ & 柔性隔热毡+碳-碳 \\
SpaceX Dragon & 25 & $\sim$20 MW/m$^{2}$ & PICA-X烧蚀 \\
火星探测器 & 30 & $\sim$5 MW/m$^{2}$ & SLA-561V烧蚀 \\
\bottomrule
\end{tabular}
\end{table}

\subsection{激波-激波相互作用}

在复杂外形（如发动机进气道唇口、控制翼面根部）周围，不同来源的激波可能相交。Edney（1968）将激波相互作用分为六种类型（I-VI型），其中IV型相互作用可产生"激波冲击加热"——局部热流放大十倍甚至数十倍，是局部热防护失效的主要原因。

激波相互作用的理论分析以激波极线图为基础，利用激波、滑流线和膨胀波之间的匹配条件建立完整解（三激波理论和自由相互作用理论）。

\begin{remark}
Edney IV型激波相互作用导致的热流峰值是高超音速飞行器热防护局部加强设计的核心输入。SpaceX的Starship不锈钢壳体的简单外形设计，初衷之一就是尽量减少激波相互作用引起的局部热斑。
\end{remark}

\section{可压缩流动的数值挑战}

虽然计算流体力学的系统论述安排在第十和第十一章，但在可压缩流动的语境下需特别指出其数值难点：

\begin{itemize}
  \item \textbf{激波捕捉。}激波是数学上的间断。中心差分格式在激波附近产生伪振荡（Gibbs现象），必须使用激波捕捉格式（TVD、ENO、WENO）或引入人工黏性。
  \item \textbf{低马赫数刚度。}当局部$Ma\approx 0$时，原始变量的Euler方程出现推进刚度——声速远大于对流速度，导致预处理（preconditioning）的必要性。
  \item \textbf{化学非平衡源项。}反应速率的高度非线性使源项刚性问题突出——需要用点隐式或分裂格式处理。
\end{itemize}

这些内容将在第十章"计算流体力学导论"和第十一章"湍流模型与数值模拟"中详细展开。

\vspace{12pt}

\noindent\textbf{本章要点回顾：}
\begin{enumerate}
  \item 声速$a = \sqrt{\gamma RT}$仅取决于温度，Mach数$Ma = U/a$是最重要的可压缩流动参数。
  \item 等熵关系给出了滞止参数与当地参数的联系，$A/A^{*}$面积比公式是Laval喷管设计的核心。
  \item 正激波由Rankine-Hugoniot关系描述，Prandtl关系$u_1 u_2 = a^{*2}$联系激波前后速度。
  \item 斜激波等效为正激波的法向分量，膨胀波由Prandtl-Meyer函数描述。
  \item 高超音速流动的特点是薄激波层、气动加热和真实气体效应，热防护是工程核心。
\end{enumerate}
